% Part: first-order-logic
% Chapter: sequent-calculus

\documentclass[../../../include/open-logic-chapter]{subfiles}

\begin{document}

\iftag{FOL}
      {\olchapter{fol}{seq}{De sequentencalculus}}
      {\olchapter{pl}{seq}{De sequentencalculus}}

\begin{editorial}
  In dit hoofdstuk gaat het over Gentzens standaardsequentencalculus LK
  voor de klassieke eerste-ordelogica. Het hoofdstuk kan nog meer
  voorbeelden en opgaven gebruiken. Met de tag ``prfLK'' bepaalt u of
  materiaal over de sequentencalculus als bewijssysteem wel of niet
  wordt opgenomen.
\end{editorial}

\olimport{rules-and-proofs}

\olimport{propositional-rules}

\iftag{FOL}{%
  \olimport{quantifier-rules}
}{}

\olimport{structural-rules}

\olimport{derivations}

\olimport{proving-things}

\iftag{FOL}{%
  \olimport{proving-things-quant}
}{}

\olimport{proof-theoretic-notions}

\olimport{provability-consistency}

\olimport{provability-propositional}

\iftag{FOL}{%
  \olimport{provability-quantifiers}
}{}

\olimport{soundness}

\iftag{FOL}{%
  \olimport{identity}
  \olimport{soundness-identity}
}{}

\OLEndChapterHook
\end{document}
